\chapter{Driven systems: perturbation theory and the rotating-wave approximation}
\label{ch:perturbations}

\section*{Chapter overview}

The Schr\"odinger equation can rarely be solved exactly when the
Hamiltonian is time-dependent. Two systematic approximations carry us
through almost every situation we will meet later. The first is
\emph{time-dependent perturbation theory}, in which a small coupling
$\op{V}(t)$ added to a known $\op{H}_0$ generates transitions whose
amplitudes are computed order by order; specialising to a continuum
of final states gives Fermi's golden rule, the standard expression
for transition rates. The second is the \emph{rotating-wave
approximation} (RWA), in which fast counter-rotating terms in the
interaction-picture Hamiltonian are dropped because they average to
zero on the timescales of interest. Both approximations are
ingredients of the Jaynes--Cummings model
(\cref{ch:jaynes-cummings}), of dispersive readout
(\cref{sec:dispersive}), and of every drive-Hamiltonian calculation
in cQED.

% --------------------------------------------------------------
\section{Time-dependent perturbation theory}
\label{sec:tdpt}

\begin{phenomenon}
A weakly perturbed system in an eigenstate of $\op{H}_0$ slowly
leaks probability amplitude into other eigenstates. To leading order
the leakage rate depends on a single matrix element of the
perturbation and on a frequency-domain mismatch between drive and
transition.
\end{phenomenon}

Let $\op{H}(t)=\op{H}_0+\op{V}(t)$ with $\op{H}_0\ket{n}=E_n\ket{n}$.
Expand the state in the interaction picture as
\begin{equation}
\ket{\psi(t)}_I = \sum_n c_n(t)\,\ket{n},
\end{equation}
where $\ket{\psi}_I=e^{i\op{H}_0 t/\hbar}\ket{\psi}_S$. Plugging into
the interaction-picture Schr\"odinger equation
$i\hbar\,\partial_t\ket{\psi}_I=\op{V}_I(t)\ket{\psi}_I$ and projecting
onto $\bra{f}$ gives
\begin{equation}
\label{eq:cdot-exact}
i\hbar\,\dot c_f(t) \;=\; \sum_n e^{i\omega_{fn}t}\,V_{fn}(t)\,c_n(t),
\qquad \omega_{fn}=(E_f-E_n)/\hbar,
\end{equation}
with $V_{fn}(t)=\bra{f}\op{V}(t)\ket{n}$.

\begin{statement}[First-order amplitude]
\label{stmt:first-order}
Starting from $\ket{i}$ at $t=0$ and assuming $\op{V}$ is small, the
amplitude to be in $\ket{f}\neq\ket{i}$ at time $t$ is
\begin{equation}
\label{eq:c-first-order}
c_f^{(1)}(t) \;=\; -\frac{i}{\hbar}\int_0^t e^{i\omega_{fi}t'}\,V_{fi}(t')\,\dd t'.
\end{equation}
\end{statement}

\begin{proof}
Set $c_n(t)=\delta_{ni}+c_n^{(1)}(t)+\ldots$ and substitute into
\cref{eq:cdot-exact}. Keeping only zeroth-order $c_n$ on the right
and $\dot c_f^{(1)}$ on the left gives
$i\hbar\,\dot c_f^{(1)}=e^{i\omega_{fi}t}V_{fi}(t)$; integrating from
$0$ to $t$ yields \cref{eq:c-first-order}.
\end{proof}

\paragraph{Transition.}
\Cref{eq:c-first-order} is exact at first order but contains only
amplitudes. To extract experimentally observable transition
\emph{rates} into a continuum of final states, we take the
long-time, dense-spectrum limit, which produces Fermi's golden rule.

% --------------------------------------------------------------
\section{Fermi's golden rule}
\label{sec:fermi}

\begin{theorem}[Fermi's golden rule]
\label{thm:fermi}
Let $\op{V}(t)=\op{V}\,e^{-i\omega t}+\op{V}^\dagger e^{+i\omega t}$ be a
harmonic perturbation, and let the final states form a continuum
with density of states $\rho_f(E)$ near $E=E_i+\hbar\omega$. The
transition rate from $\ket{i}$ to states near energy
$E_f=E_i+\hbar\omega$ is
\begin{equation}
\label{eq:fermi}
\Gamma_{i\to f} \;=\; \frac{2\pi}{\hbar}\,\bigl|\bra{f}\op{V}\ket{i}\bigr|^2\,\rho_f(E_f).
\end{equation}
\end{theorem}

\begin{proof}
From \cref{eq:c-first-order} with $V_{fi}(t)=V_{fi}\,e^{-i\omega t}$,
\begin{equation}
c_f^{(1)}(t) = -\frac{i}{\hbar}V_{fi}\,
\frac{e^{i(\omega_{fi}-\omega)t}-1}{i(\omega_{fi}-\omega)},
\end{equation}
so
\begin{equation}
\bigl|c_f^{(1)}(t)\bigr|^2
= \frac{|V_{fi}|^2}{\hbar^2}\,\frac{4\sin^2[(\omega_{fi}-\omega)t/2]}{(\omega_{fi}-\omega)^2}.
\end{equation}
For large $t$ the function on the right is sharply peaked at
$\omega_{fi}=\omega$ with integral $2\pi t$; in the distributional
sense $4\sin^2(xt/2)/x^2\to 2\pi t\,\delta(x)$. Summing over a
continuum of final states with density $\rho_f(E_f)$ and dividing by
$t$ gives \cref{eq:fermi}.
\end{proof}

\begin{example}[Spontaneous emission and the Purcell rate]
The decay rate of an excited atom into the continuum of vacuum
electromagnetic modes follows from Fermi's golden rule with
$|V_{fi}|^2$ proportional to the dipole matrix element squared and
$\rho_f(E_f)$ to the photon density of states. In a cavity,
modifying $\rho_f$ at the atomic frequency accelerates or suppresses
the decay --- the Purcell effect of \cref{sec:purcell}.
\end{example}

\paragraph{Transition.}
Fermi's golden rule applies to weak perturbations and incoherent
transitions. For \emph{strong} but \emph{near-resonant} drives, like
the microwave tones used to control qubits, a different
approximation is more useful: drop the fast off-resonant terms in
the interaction-picture Hamiltonian.

% --------------------------------------------------------------
\section{The rotating-wave approximation}
\label{sec:rwa}

\begin{phenomenon}
A near-resonant drive on a quantum system produces both a slow,
``co-rotating'' amplitude that drives transitions and a fast,
``counter-rotating'' one that averages to zero on the timescale of
the slow dynamics. Provided the slow and fast scales are well
separated, the counter-rotating part can be discarded.
\end{phenomenon}

\begin{statement}[RWA, generic form]
\label{stmt:rwa}
Consider $\op{H}(t)=\op{H}_0+\op{V}(t)$ with
$\op{H}_0\ket{n}=E_n\ket{n}$ and $\op{V}(t)$ harmonic at frequency
$\omega$. In the interaction picture,
$\op{V}_I(t)=\sum_{m,n} V_{mn}\,e^{i\omega_{mn}t}\,e^{\pm i\omega t}\,
\ket{m}\bra{n}$. Terms with $|\omega_{mn}\pm\omega|\gg|V_{mn}|/\hbar$
oscillate fast on the scale of $\hbar/V$ and integrate to a
negligible contribution; the RWA consists in dropping them and
keeping only the resonant terms with $\omega_{mn}\pm\omega\approx0$.
\end{statement}

\paragraph{Derivation, two-level example.}
For the driven TLS of \cref{sec:rabi} with
$\op{V}(t)=\Omega\cos(\omega_d t)\sigma_x$, expand
$\sigma_x=\sigma_++\sigma_-$ and write
\begin{equation}
\op{V}_I(t) = \tfrac{\Omega}{2}\bigl(\sigma_+ e^{i\omega_q t}+\sigma_- e^{-i\omega_q t}\bigr)
\bigl(e^{i\omega_d t}+e^{-i\omega_d t}\bigr).
\end{equation}
Multiplying out, four terms appear: two oscillate at
$\omega_q\pm\omega_d$ (slow, when $\omega_d\approx\omega_q$) and two
at $\omega_q+\omega_d$ or $-(\omega_q+\omega_d)$ (always fast).
Dropping the latter,
\begin{equation}
\op{V}_I^\text{RWA}(t)
= \tfrac{\Omega}{2}\bigl(\sigma_+ e^{+i(\omega_q-\omega_d)t}+\sigma_- e^{-i(\omega_q-\omega_d)t}\bigr).
\end{equation}
Going one step further into the frame rotating at the drive
frequency, $\op{U}_d(t)=e^{-i\omega_d t\sigma_z/2}$, removes the
remaining time dependence and reproduces \cref{eq:rabi-h-rwa}, the
static Hamiltonian whose dynamics gave the Rabi formula.

\begin{remark}[Validity]
The RWA is exact only in the limit
$\Omega\ll \omega_q+\omega_d$ and $|\omega_q-\omega_d|\ll\omega_q+\omega_d$.
At very strong driving the kept terms produce dynamics on a timescale
comparable to the discarded fast oscillations, and corrections such
as the Bloch--Siegert shift become measurable.
\end{remark}

\begin{example}[Why we needed it earlier]
The Rabi formula in \cref{thm:rabi} was derived using exactly this
approximation. Every drive Hamiltonian we will meet from
\cref{ch:jaynes-cummings} onwards --- the Jaynes--Cummings coupling,
the dispersive readout drive, the cross-resonance two-qubit gate ---
takes its standard textbook form only after invoking the RWA.
\end{example}

% --------------------------------------------------------------
\section*{Chapter summary and outlook}

This chapter completed the toolkit needed to handle time-dependent
problems: time-dependent perturbation theory gives transition
amplitudes order by order, Fermi's golden rule converts those
amplitudes into rates whenever a continuum of final states is
available, and the rotating-wave approximation simplifies
near-resonant drives to a static Hamiltonian in the appropriate
rotating frame.

Together with the postulates and operators of
\cref{ch:postulates}, the model systems of \cref{ch:model-systems},
and the composite-system formalism of \cref{ch:composite}, we have
everything we need to start doing physics. In Part~II we will use
this machinery to describe systems coupled to a noisy environment ---
producing the Lindblad master equation, the relaxation and dephasing
times $T_1, T_2$, and the input--output formalism. From there
Part~III turns to quantum optics proper, and Part~IV to
superconducting circuits, where every approximation introduced in
this chapter will earn its keep.
